\worksheettitle
  {M04\(\star\). Два положения точки}
  {\modulemark{M04\(\star\)} Усложнение}

\studentnameblock

\begin{notebox}[Главная идея]
  От заданной точки расстояние можно отложить влево и вправо. Каждый
  полученный вариант нужно проверить по условию и границам отрезка.
\end{notebox}

\twopositionsfigure

\begin{practicebox}[1. Рассмотрите два случая]
  $AB=10$ см, $AC=6$ см, $CD=2$ см. Отметьте оба допустимых положения
  точки $D$ и найдите $BD$ в каждом случае.

  \studenttwocasesfigure

  Ответы: $BD=\answerblank[20mm]$ или $BD=\answerblank[20mm]$.
\end{practicebox}

\begin{practicebox}[2. Отберите допустимые случаи]
  Точка $D$ должна лежать на отрезке $AB=12$ см.
  \begin{enumerate}[label=\alph*)]
    \item $AC=5$ см, $CD=4$ см: число положений \answerblank[15mm];
    \item $AC=2$ см, $CD=5$ см: число положений \answerblank[15mm];
    \item $AC=3$ см, $CD=13$ см: число положений \answerblank[15mm].
  \end{enumerate}
\end{practicebox}

\begin{checkpointbox}[3. Докажите полноту]
  Для $AB=15$ см, $AC=7$ см, $CD=6$ см найдите все значения $BD$.
  Объясните, почему других ответов нет. \writinglines{4}
\end{checkpointbox}
